%------------------------------------------------------------------------------%
\documentclass[fleqn,utf8,aspectratio=169,12pt]{beamer}

\usepackage{gaal}

\title{Espaços Vetoriais}

%------------------------------------------------------------------------------%
\begin{document}

\addtitle

%------------------------------------------------------------------------------%
\section{Espaços Vetoriais}

%------------------------------------------------------------------------------%
\begin{frame}{Motivação}
  \begin{itemize}[<+->]
    \setlength{\itemsep}{1em}
    \item Vetores aparecem em diferentes contextos
    \item As operações podem ter naturezas diferentes
    \item A mesma estrutura algébrica está presente em todos eles
  \end{itemize}
\end{frame}

%------------------------------------------------------------------------------%
\begin{frame}{Abstração}
  O vetor não precisa ser uma seta

  \pause
  Qualquer elemento matemático que podemos
  \pause
  \begin{enumerate}
    \item somar
    \pause
    \item multiplicar por escalar
  \end{enumerate}
  \pause
  Pode ser pensado como vetor
\end{frame}

%------------------------------------------------------------------------------%
\begin{frame}{Corpo}
  Um conjunto $K$ que possui operações de soma e produto

  \pause
  \medskip
  com as mesmas propriedades da soma e produto dos números reais

  \pause
  \medskip
  é chamado de \emph{Corpo}
\end{frame}

%------------------------------------------------------------------------------%
\begin{frame}{Espaços Vetoriais}
  Um \emph{Espaço vetorial} consiste de

  \pause
  \quad um conjunto \quad $V$

  \pause
  \quad um corpo \quad\quad $K$

  \pause
  \quad uma operação de soma \quad \(+\colon V\times V \to V\)

  \pause
  \quad um produto dor escalar \quad \(\cdot\,\colon V\times K \to V\)
\end{frame}

%------------------------------------------------------------------------------%
\begin{frame}{Operações}
  Para quais quer \(\vec{v}\), \(\vec{w}\in V\) e \(\alpha\in K\)

  \pause
  \[
    \vec{v}+\vec{w} \in V
  \]

  \pause
  \[
    \alpha\vec{v} \in V
  \]
\end{frame}

%------------------------------------------------------------------------------%
\begin{frame}{Axiomas}
  Para quaisquer \(\vec{u}\), \(\vec{v}\), \(\vec{w}\in V\)
  \pause
  e \(\alpha\), \(\beta\in K\)
  \[
    \pause
    \vec{u}+\vec{v} \;=\; \vec{v}+\vec{u}
  \]

  \[
    \pause
    (\vec{u}+\vec{v})+\vec{w} \;=\; \vec{u}+(\vec{v}+\vec{w})
  \]

  \[
    \pause
    \exists\,\vec{0}\in V
    \qquad
    \vec{v}+\vec{0} = \vec{v}
  \]

  \[
    \pause
    \forall\,\vec{v}\in V
    \qquad
    \exists\,(-\vec{v})\in V
    \qquad
    \vec{v}+(-\vec{v}) = \vec{0}
  \]
\end{frame}

%------------------------------------------------------------------------------%
\begin{frame}{Axiomas}
  Para quaisquer \(\vec{u}\), \(\vec{v}\in V\)
  \pause
  e \(\alpha\), \(\beta\in K\)
  \[
    \pause
    \alpha(\vec{u}+\vec{v}) \;=\; \alpha\vec{u}+\alpha\vec{v}
  \]

  \[
    \pause
    (\alpha+\beta)\vec{v} \;=\; \alpha\vec{v}+\beta\vec{v}
  \]

  \[
    \pause
    \alpha(\beta\vec{v}) \;=\; (\alpha\beta)\vec{v}
  \]

  \[
    \pause
    1\vec{v} \;=\; \vec{v}
  \]
\end{frame}

%------------------------------------------------------------------------------%
\section{Espaços Vetoriais Usuais}

%------------------------------------------------------------------------------%
\begin{frame}{O espaço $\R^2$}
  Conjunto
  \[
    \R^2 = \left\{ (x, y) \mid x, y \in\R \right\}
  \]

  \pause
  Soma
  \[
    (x_1, y_1) + (x_2, y_2) = (x_1+x_2,,\ y_1+y_2)
  \]

  \pause
  Produto por escalar
  \[
    \alpha(x, y) = (\alpha x, \alpha y)
  \]
\end{frame}

%------------------------------------------------------------------------------%
\begin{frame}{O espaço $\R^n$}
  Conjunto
  \[
    \R^n = \left\{ (x_1, x_2, \ldots, x_n) \mid x_i\in\R \right\}
  \]

  \pause
  Soma
  \[
    (x_1, \ldots, x_n) + (y_1, \ldots, y_n) = (x_1 + y_1, \ldots, x_n + y_n)
  \]

  \pause
  Produto por escalar
  \[
    \alpha(x_1, \ldots, x_n) = (\alpha x_1, \ldots, \alpha x_n)
  \]
\end{frame}

%------------------------------------------------------------------------------%
\begin{frame}{Polinômios de Grau $n$}
  Conjunto
  \[
    \mathcal{P}_n = \left\{ a_0 + a_1x + \cdots + a_nx^n \mid a_i\in\R \right\}
  \]

  \pause
  Soma
  \[
    p(x) + q(x)
  \]

  \pause
  Produto por escalar
  \[
    \alpha p(x)
  \]
\end{frame}

%------------------------------------------------------------------------------%
\begin{frame}{Funções Reais}
  Conjunto de todas as funções reais de variável real
  \[
    \mathcal{F}\left(\R, \R\right)
  \]
  \[
    f\colon \R \to \R
  \]

  \pause
  Soma
  \[
    (f+g)(x) = f(x) + g(x)
  \]

  \pause
  Produto por escalar
  \[
    (\alpha f)(x) = \alpha f(x)
  \]
\end{frame}

%------------------------------------------------------------------------------%
\begin{frame}{Matrizes}
\begin{columns}
\column{0.4\textwidth}
  Conjunto das matrizes reais
  \[
    M_{m\times n}\left(\R\right)
  \]
  \pause
  \[
    A =
    \begin{pmatrix}
      a_{11} & \cdots & a_{1n} \\
      \vdots & \ddots & \vdots \\
      a_{m1} & \cdots & a_{mn}
    \end{pmatrix}
  \]
\column{0.4\textwidth}
  \pause
  Soma
  \[
    A + B
    = \left( a_{ij} + b_{ij} \right)
  \]

  \pause
  Produto por escalar
  \[
    \alpha A = \left(\alpha a_{ij}\right)
  \]
\end{columns}
\end{frame}

%------------------------------------------------------------------------------%
\begin{frame}{Nem Sempre o Conjunto é um Espaço Vetorial}
  O conjunto das matrizes
  \(
    M_{n\times n}(\R)
  \)
  é um espaço vetorial

  \pause
  Mas o conjunto das \emph{matrizes invertíveis}
  \[
    GL_n(\R)
  \]
  \pause
  não é um espaço vetorial com as operações usuais

  \[
    \pause
    I \in GL_n(\R)
    \pause
    \quad\text{ mas }\quad
    0I \notin GL_n(\R)
  \]
\end{frame}

%%------------------------------------------------------------------------------%
%\section{Lista Mínima}
%
%%------------------------------------------------------------------------------%
%\begin{frame}{Lista Mínima}
%  \begin{enumerate}
%    \item
%  \end{enumerate}
%
%  \vfill
%  Atenção: A prova é baseada no livro, não nas apresentações
%\end{frame}

%------------------------------------------------------------------------------%
\end{document}
%------------------------------------------------------------------------------%
